\section{嘉道危机其余账本事件锚点}
本节补齐嘉道危机余下事件的导航入口，确保财政、边贸、河工、战争、口岸与地方社会材料均可追溯。
\subsection{1790—1799}
\eventanchor{EQ-1798-MILITIA-MOBILIZATION-DECISION}{围绕「朝廷加强对乡勇、团练和…}嘉庆三年（1798年），围绕「朝廷加强对乡勇、团练和地方绅士资源的动员以应对白莲教战争」的命令、调度、照会或呈报形成独立的中央—地方行动文书链。核心取证：《清仁宗实录》嘉庆三年相关条；地方志、督抚奏折及地方档案。该项锚定的是命令或决策环节，命令抵达和结果仍须另外核验。来源保留：奏折、上谕、部议、军机档或条约可证明行政动作和机构立场，不自动证明所有地区、群体或对手按其意图行动。
\eventanchor{EQ-1798-MILITIA-MOBILIZATION-AFTERMATH}{「朝廷加强对乡勇、团练和地方…}嘉庆三年（1798年），「朝廷加强对乡勇、团练和地方绅士资源的动员以应对白莲教战争」引发的善后、执行或地方回应构成可由不同材料追踪的后续节点。核心取证：《清仁宗实录》嘉庆三年相关条；地方志、督抚奏折及地方档案；相关地区的档案、志书、契约、报刊或对方材料。该项锚定的是后续追踪问题，影响范围须按地区与群体分别判断。来源保留：战后、改革后或条约后的记忆、地方记录和官方总结常具有事后选择性；后续影响须按地点、群体和时间分层，不作整体化推断。
\eventanchor{EQ-1799-HESHEN-PURGE-PRELUDE}{「乾隆帝去世后，嘉庆帝逮治和…}嘉庆四年（1799年），「乾隆帝去世后，嘉庆帝逮治和珅并清理其政治网络」之前的条件、信息和争议已通过中央与地方材料显现，构成该事件可追踪的前史节点。核心取证：《清仁宗实录》嘉庆四年相关条；宫中朱批奏折、内阁大库题本与《清史稿》相关志传；同时期地方档案、地方志、日记或对方材料应按具体地区补检。该项锚定的是前史条件与信息背景，不把条件直接写成唯一因果。来源保留：官修编年和地方材料的记录密度与立场并不相同；本行不将条件、传闻、呈报或后见解释直接等同为确定因果。
\eventanchor{EQ-1799-HESHEN-PURGE-DECISION}{围绕「乾隆帝去世后，嘉庆帝逮…}嘉庆四年（1799年），围绕「乾隆帝去世后，嘉庆帝逮治和珅并清理其政治网络」的命令、调度、照会或呈报形成独立的中央—地方行动文书链。核心取证：《清仁宗实录》嘉庆四年相关条；宫中朱批奏折、内阁大库题本与《清史稿》相关志传。该项锚定的是命令或决策环节，命令抵达和结果仍须另外核验。来源保留：奏折、上谕、部议、军机档或条约可证明行政动作和机构立场，不自动证明所有地区、群体或对手按其意图行动。
\eventanchor{EQ-1799-HESHEN-PURGE-AFTERMATH}{「乾隆帝去世后，嘉庆帝逮治和…}嘉庆四年（1799年），「乾隆帝去世后，嘉庆帝逮治和珅并清理其政治网络」引发的善后、执行或地方回应构成可由不同材料追踪的后续节点。核心取证：《清仁宗实录》嘉庆四年相关条；宫中朱批奏折、内阁大库题本与《清史稿》相关志传；相关地区的档案、志书、契约、报刊或对方材料。该项锚定的是后续追踪问题，影响范围须按地区与群体分别判断。来源保留：战后、改革后或条约后的记忆、地方记录和官方总结常具有事后选择性；后续影响须按地点、群体和时间分层，不作整体化推断。
\subsection{1800—1809}
\eventanchor{EQ-1801-WHITE-LOTUS-SUPPRESSION-PRELUDE}{「清军在若干战区逐步压缩白莲…}嘉庆六年（1801年），「清军在若干战区逐步压缩白莲教起事者活动空间」之前的条件、信息和争议已通过中央与地方材料显现，构成该事件可追踪的前史节点。核心取证：《清仁宗实录》嘉庆六年相关条；军机处档、战事方略及地方奏报；同时期地方档案、地方志、日记或对方材料应按具体地区补检。该项锚定的是前史条件与信息背景，不把条件直接写成唯一因果。来源保留：官修编年和地方材料的记录密度与立场并不相同；本行不将条件、传闻、呈报或后见解释直接等同为确定因果。
\eventanchor{EQ-1801-WHITE-LOTUS-SUPPRESSION-DECISION}{围绕「清军在若干战区逐步压缩…}嘉庆六年（1801年），围绕「清军在若干战区逐步压缩白莲教起事者活动空间」的命令、调度、照会或呈报形成独立的中央—地方行动文书链。核心取证：《清仁宗实录》嘉庆六年相关条；军机处档、战事方略及地方奏报。该项锚定的是命令或决策环节，命令抵达和结果仍须另外核验。来源保留：奏折、上谕、部议、军机档或条约可证明行政动作和机构立场，不自动证明所有地区、群体或对手按其意图行动。
\eventanchor{EQ-1801-WHITE-LOTUS-SUPPRESSION-AFTERMATH}{「清军在若干战区逐步压缩白莲…}嘉庆六年（1801年），「清军在若干战区逐步压缩白莲教起事者活动空间」引发的善后、执行或地方回应构成可由不同材料追踪的后续节点。核心取证：《清仁宗实录》嘉庆六年相关条；军机处档、战事方略及地方奏报；相关地区的档案、志书、契约、报刊或对方材料。该项锚定的是后续追踪问题，影响范围须按地区与群体分别判断。来源保留：战后、改革后或条约后的记忆、地方记录和官方总结常具有事后选择性；后续影响须按地点、群体和时间分层，不作整体化推断。
\eventanchor{EQ-1802-WAR-REFUGEE-RELIEF-PRELUDE}{「战区流民、复垦和赈济成为地…}嘉庆七年（1802年），「战区流民、复垦和赈济成为地方善后事务」之前的条件、信息和争议已通过中央与地方材料显现，构成该事件可追踪的前史节点。核心取证：《清仁宗实录》嘉庆七年相关条；地方志、督抚奏折及地方档案；同时期地方档案、地方志、日记或对方材料应按具体地区补检。该项锚定的是前史条件与信息背景，不把条件直接写成唯一因果。来源保留：官修编年和地方材料的记录密度与立场并不相同；本行不将条件、传闻、呈报或后见解释直接等同为确定因果。
\eventanchor{EQ-1802-WAR-REFUGEE-RELIEF-DECISION}{围绕「战区流民、复垦和赈济成…}嘉庆七年（1802年），围绕「战区流民、复垦和赈济成为地方善后事务」的命令、调度、照会或呈报形成独立的中央—地方行动文书链。核心取证：《清仁宗实录》嘉庆七年相关条；地方志、督抚奏折及地方档案。该项锚定的是命令或决策环节，命令抵达和结果仍须另外核验。来源保留：奏折、上谕、部议、军机档或条约可证明行政动作和机构立场，不自动证明所有地区、群体或对手按其意图行动。
\eventanchor{EQ-1802-WAR-REFUGEE-RELIEF-AFTERMATH}{「战区流民、复垦和赈济成为地…}嘉庆七年（1802年），「战区流民、复垦和赈济成为地方善后事务」引发的善后、执行或地方回应构成可由不同材料追踪的后续节点。核心取证：《清仁宗实录》嘉庆七年相关条；地方志、督抚奏折及地方档案；相关地区的档案、志书、契约、报刊或对方材料。该项锚定的是后续追踪问题，影响范围须按地区与群体分别判断。来源保留：战后、改革后或条约后的记忆、地方记录和官方总结常具有事后选择性；后续影响须按地点、群体和时间分层，不作整体化推断。
\eventanchor{EQ-1803-FORBIDDEN-CITY-ATTACK-PRELUDE}{「陈德等闯入紫禁城，暴露京师…}嘉庆八年（1803年），「陈德等闯入紫禁城，暴露京师警卫与内廷管理问题」之前的条件、信息和争议已通过中央与地方材料显现，构成该事件可追踪的前史节点。核心取证：《清仁宗实录》嘉庆八年相关条；宫中朱批奏折、内阁大库题本与《清史稿》相关志传；同时期地方档案、地方志、日记或对方材料应按具体地区补检。该项锚定的是前史条件与信息背景，不把条件直接写成唯一因果。来源保留：官修编年和地方材料的记录密度与立场并不相同；本行不将条件、传闻、呈报或后见解释直接等同为确定因果。
\eventanchor{EQ-1803-FORBIDDEN-CITY-ATTACK-DECISION}{围绕「陈德等闯入紫禁城，暴露…}嘉庆八年（1803年），围绕「陈德等闯入紫禁城，暴露京师警卫与内廷管理问题」的命令、调度、照会或呈报形成独立的中央—地方行动文书链。核心取证：《清仁宗实录》嘉庆八年相关条；宫中朱批奏折、内阁大库题本与《清史稿》相关志传。该项锚定的是命令或决策环节，命令抵达和结果仍须另外核验。来源保留：奏折、上谕、部议、军机档或条约可证明行政动作和机构立场，不自动证明所有地区、群体或对手按其意图行动。
\eventanchor{EQ-1803-FORBIDDEN-CITY-ATTACK-AFTERMATH}{「陈德等闯入紫禁城，暴露京师…}嘉庆八年（1803年），「陈德等闯入紫禁城，暴露京师警卫与内廷管理问题」引发的善后、执行或地方回应构成可由不同材料追踪的后续节点。核心取证：《清仁宗实录》嘉庆八年相关条；宫中朱批奏折、内阁大库题本与《清史稿》相关志传；相关地区的档案、志书、契约、报刊或对方材料。该项锚定的是后续追踪问题，影响范围须按地区与群体分别判断。来源保留：战后、改革后或条约后的记忆、地方记录和官方总结常具有事后选择性；后续影响须按地点、群体和时间分层，不作整体化推断。
\eventanchor{EQ-1805-WHITE-LOTUS-WAR-END-PRELUDE}{「白莲教大规模战争基本结束，…}嘉庆十年（1805年），「白莲教大规模战争基本结束，清廷转入善后和治安重整」之前的条件、信息和争议已通过中央与地方材料显现，构成该事件可追踪的前史节点。核心取证：《清仁宗实录》嘉庆十年相关条；军机处档、战事方略及地方奏报；同时期地方档案、地方志、日记或对方材料应按具体地区补检。该项锚定的是前史条件与信息背景，不把条件直接写成唯一因果。来源保留：官修编年和地方材料的记录密度与立场并不相同；本行不将条件、传闻、呈报或后见解释直接等同为确定因果。
\eventanchor{EQ-1805-WHITE-LOTUS-WAR-END-DECISION}{围绕「白莲教大规模战争基本结…}嘉庆十年（1805年），围绕「白莲教大规模战争基本结束，清廷转入善后和治安重整」的命令、调度、照会或呈报形成独立的中央—地方行动文书链。核心取证：《清仁宗实录》嘉庆十年相关条；军机处档、战事方略及地方奏报。该项锚定的是命令或决策环节，命令抵达和结果仍须另外核验。来源保留：奏折、上谕、部议、军机档或条约可证明行政动作和机构立场，不自动证明所有地区、群体或对手按其意图行动。
\eventanchor{EQ-1805-WHITE-LOTUS-WAR-END-AFTERMATH}{「白莲教大规模战争基本结束，…}嘉庆十年（1805年），「白莲教大规模战争基本结束，清廷转入善后和治安重整」引发的善后、执行或地方回应构成可由不同材料追踪的后续节点。核心取证：《清仁宗实录》嘉庆十年相关条；军机处档、战事方略及地方奏报；相关地区的档案、志书、契约、报刊或对方材料。该项锚定的是后续追踪问题，影响范围须按地区与群体分别判断。来源保留：战后、改革后或条约后的记忆、地方记录和官方总结常具有事后选择性；后续影响须按地点、群体和时间分层，不作整体化推断。
\eventanchor{EQ-1806-TSAI-QIAN-PIRACY-PRELUDE}{「蔡牵等海盗集团持续活动，清…}嘉庆十一年（1806年），「蔡牵等海盗集团持续活动，清廷调整水师与沿海防务」之前的条件、信息和争议已通过中央与地方材料显现，构成该事件可追踪的前史节点。核心取证：《清仁宗实录》嘉庆十一年相关条；军机处档、战事方略及地方奏报；同时期地方档案、地方志、日记或对方材料应按具体地区补检。该项锚定的是前史条件与信息背景，不把条件直接写成唯一因果。来源保留：官修编年和地方材料的记录密度与立场并不相同；本行不将条件、传闻、呈报或后见解释直接等同为确定因果。
\eventanchor{EQ-1806-TSAI-QIAN-PIRACY-DECISION}{围绕「蔡牵等海盗集团持续活动…}嘉庆十一年（1806年），围绕「蔡牵等海盗集团持续活动，清廷调整水师与沿海防务」的命令、调度、照会或呈报形成独立的中央—地方行动文书链。核心取证：《清仁宗实录》嘉庆十一年相关条；军机处档、战事方略及地方奏报。该项锚定的是命令或决策环节，命令抵达和结果仍须另外核验。来源保留：奏折、上谕、部议、军机档或条约可证明行政动作和机构立场，不自动证明所有地区、群体或对手按其意图行动。
\eventanchor{EQ-1806-TSAI-QIAN-PIRACY-AFTERMATH}{「蔡牵等海盗集团持续活动，清…}嘉庆十一年（1806年），「蔡牵等海盗集团持续活动，清廷调整水师与沿海防务」引发的善后、执行或地方回应构成可由不同材料追踪的后续节点。核心取证：《清仁宗实录》嘉庆十一年相关条；军机处档、战事方略及地方奏报；相关地区的档案、志书、契约、报刊或对方材料。该项锚定的是后续追踪问题，影响范围须按地区与群体分别判断。来源保留：战后、改革后或条约后的记忆、地方记录和官方总结常具有事后选择性；后续影响须按地点、群体和时间分层，不作整体化推断。
\eventanchor{EQ-1807-COASTAL-NAVAL-REFORM-PRELUDE}{「清廷继续通过水师调度、招抚…}嘉庆十二年（1807年），「清廷继续通过水师调度、招抚和地方协作应对海盗」之前的条件、信息和争议已通过中央与地方材料显现，构成该事件可追踪的前史节点。核心取证：《清仁宗实录》嘉庆十二年相关条；宫中朱批奏折、内阁大库题本与《清史稿》相关志传；同时期地方档案、地方志、日记或对方材料应按具体地区补检。该项锚定的是前史条件与信息背景，不把条件直接写成唯一因果。来源保留：官修编年和地方材料的记录密度与立场并不相同；本行不将条件、传闻、呈报或后见解释直接等同为确定因果。
\eventanchor{EQ-1807-COASTAL-NAVAL-REFORM-DECISION}{围绕「清廷继续通过水师调度、…}嘉庆十二年（1807年），围绕「清廷继续通过水师调度、招抚和地方协作应对海盗」的命令、调度、照会或呈报形成独立的中央—地方行动文书链。核心取证：《清仁宗实录》嘉庆十二年相关条；宫中朱批奏折、内阁大库题本与《清史稿》相关志传。该项锚定的是命令或决策环节，命令抵达和结果仍须另外核验。来源保留：奏折、上谕、部议、军机档或条约可证明行政动作和机构立场，不自动证明所有地区、群体或对手按其意图行动。
\eventanchor{EQ-1807-COASTAL-NAVAL-REFORM-AFTERMATH}{「清廷继续通过水师调度、招抚…}嘉庆十二年（1807年），「清廷继续通过水师调度、招抚和地方协作应对海盗」引发的善后、执行或地方回应构成可由不同材料追踪的后续节点。核心取证：《清仁宗实录》嘉庆十二年相关条；宫中朱批奏折、内阁大库题本与《清史稿》相关志传；相关地区的档案、志书、契约、报刊或对方材料。该项锚定的是后续追踪问题，影响范围须按地区与群体分别判断。来源保留：战后、改革后或条约后的记忆、地方记录和官方总结常具有事后选择性；后续影响须按地点、群体和时间分层，不作整体化推断。
\eventanchor{EQ-1809-SOUTH-CHINA-SEA-CAMPAIGN-PRELUDE}{「清军、地方势力与外来船只在…}嘉庆十四年（1809年），「清军、地方势力与外来船只在华南海域共同影响剿盗行动」之前的条件、信息和争议已通过中央与地方材料显现，构成该事件可追踪的前史节点。核心取证：《清仁宗实录》嘉庆十四年相关条；军机处档、战事方略及地方奏报；同时期地方档案、地方志、日记或对方材料应按具体地区补检。该项锚定的是前史条件与信息背景，不把条件直接写成唯一因果。来源保留：官修编年和地方材料的记录密度与立场并不相同；本行不将条件、传闻、呈报或后见解释直接等同为确定因果。
\eventanchor{EQ-1809-SOUTH-CHINA-SEA-CAMPAIGN-DECISION}{围绕「清军、地方势力与外来船…}嘉庆十四年（1809年），围绕「清军、地方势力与外来船只在华南海域共同影响剿盗行动」的命令、调度、照会或呈报形成独立的中央—地方行动文书链。核心取证：《清仁宗实录》嘉庆十四年相关条；军机处档、战事方略及地方奏报。该项锚定的是命令或决策环节，命令抵达和结果仍须另外核验。来源保留：奏折、上谕、部议、军机档或条约可证明行政动作和机构立场，不自动证明所有地区、群体或对手按其意图行动。
\eventanchor{EQ-1809-SOUTH-CHINA-SEA-CAMPAIGN-AFTERMATH}{「清军、地方势力与外来船只在…}嘉庆十四年（1809年），「清军、地方势力与外来船只在华南海域共同影响剿盗行动」引发的善后、执行或地方回应构成可由不同材料追踪的后续节点。核心取证：《清仁宗实录》嘉庆十四年相关条；军机处档、战事方略及地方奏报；相关地区的档案、志书、契约、报刊或对方材料。该项锚定的是后续追踪问题，影响范围须按地区与群体分别判断。来源保留：战后、改革后或条约后的记忆、地方记录和官方总结常具有事后选择性；后续影响须按地点、群体和时间分层，不作整体化推断。
\subsection{1810—1819}
\eventanchor{EQ-1810-ZHANG-BAOZAI-SURRENDER-PRELUDE}{「张保仔接受招安，华南大海盗…}嘉庆十五年（1810年），「张保仔接受招安，华南大海盗联盟瓦解」之前的条件、信息和争议已通过中央与地方材料显现，构成该事件可追踪的前史节点。核心取证：《清仁宗实录》嘉庆十五年相关条；军机处档、战事方略及地方奏报；同时期地方档案、地方志、日记或对方材料应按具体地区补检。该项锚定的是前史条件与信息背景，不把条件直接写成唯一因果。来源保留：官修编年和地方材料的记录密度与立场并不相同；本行不将条件、传闻、呈报或后见解释直接等同为确定因果。
\eventanchor{EQ-1810-ZHANG-BAOZAI-SURRENDER-DECISION}{围绕「张保仔接受招安，华南大…}嘉庆十五年（1810年），围绕「张保仔接受招安，华南大海盗联盟瓦解」的命令、调度、照会或呈报形成独立的中央—地方行动文书链。核心取证：《清仁宗实录》嘉庆十五年相关条；军机处档、战事方略及地方奏报。该项锚定的是命令或决策环节，命令抵达和结果仍须另外核验。来源保留：奏折、上谕、部议、军机档或条约可证明行政动作和机构立场，不自动证明所有地区、群体或对手按其意图行动。
\eventanchor{EQ-1810-ZHANG-BAOZAI-SURRENDER-AFTERMATH}{「张保仔接受招安，华南大海盗…}嘉庆十五年（1810年），「张保仔接受招安，华南大海盗联盟瓦解」引发的善后、执行或地方回应构成可由不同材料追踪的后续节点。核心取证：《清仁宗实录》嘉庆十五年相关条；军机处档、战事方略及地方奏报；相关地区的档案、志书、契约、报刊或对方材料。该项锚定的是后续追踪问题，影响范围须按地区与群体分别判断。来源保留：战后、改革后或条约后的记忆、地方记录和官方总结常具有事后选择性；后续影响须按地点、群体和时间分层，不作整体化推断。
\eventanchor{EQ-1811-BANNER-LIVELIHOOD-PRELUDE}{「八旗生计、债务和人口供养问…}嘉庆十六年（1811年），「八旗生计、债务和人口供养问题持续成为中央与地方的治理议题」之前的条件、信息和争议已通过中央与地方材料显现，构成该事件可追踪的前史节点。核心取证：《清仁宗实录》嘉庆十六年相关条；地方志、督抚奏折及地方档案；同时期地方档案、地方志、日记或对方材料应按具体地区补检。该项锚定的是前史条件与信息背景，不把条件直接写成唯一因果。来源保留：官修编年和地方材料的记录密度与立场并不相同；本行不将条件、传闻、呈报或后见解释直接等同为确定因果。
\eventanchor{EQ-1811-BANNER-LIVELIHOOD-DECISION}{围绕「八旗生计、债务和人口供…}嘉庆十六年（1811年），围绕「八旗生计、债务和人口供养问题持续成为中央与地方的治理议题」的命令、调度、照会或呈报形成独立的中央—地方行动文书链。核心取证：《清仁宗实录》嘉庆十六年相关条；地方志、督抚奏折及地方档案。该项锚定的是命令或决策环节，命令抵达和结果仍须另外核验。来源保留：奏折、上谕、部议、军机档或条约可证明行政动作和机构立场，不自动证明所有地区、群体或对手按其意图行动。
\eventanchor{EQ-1811-BANNER-LIVELIHOOD-AFTERMATH}{「八旗生计、债务和人口供养问…}嘉庆十六年（1811年），「八旗生计、债务和人口供养问题持续成为中央与地方的治理议题」引发的善后、执行或地方回应构成可由不同材料追踪的后续节点。核心取证：《清仁宗实录》嘉庆十六年相关条；地方志、督抚奏折及地方档案；相关地区的档案、志书、契约、报刊或对方材料。该项锚定的是后续追踪问题，影响范围须按地区与群体分别判断。来源保留：战后、改革后或条约后的记忆、地方记录和官方总结常具有事后选择性；后续影响须按地点、群体和时间分层，不作整体化推断。
\eventanchor{EQ-1813-TIANLI-UPRISING-PRELUDE}{「天理教起事者攻入紫禁城部分…}嘉庆十八年（1813年），「天理教起事者攻入紫禁城部分区域，京师发生严重安全危机」之前的条件、信息和争议已通过中央与地方材料显现，构成该事件可追踪的前史节点。核心取证：《清仁宗实录》嘉庆十八年相关条；军机处档、战事方略及地方奏报；同时期地方档案、地方志、日记或对方材料应按具体地区补检。该项锚定的是前史条件与信息背景，不把条件直接写成唯一因果。来源保留：官修编年和地方材料的记录密度与立场并不相同；本行不将条件、传闻、呈报或后见解释直接等同为确定因果。
\eventanchor{EQ-1813-TIANLI-UPRISING-DECISION}{围绕「天理教起事者攻入紫禁城…}嘉庆十八年（1813年），围绕「天理教起事者攻入紫禁城部分区域，京师发生严重安全危机」的命令、调度、照会或呈报形成独立的中央—地方行动文书链。核心取证：《清仁宗实录》嘉庆十八年相关条；军机处档、战事方略及地方奏报。该项锚定的是命令或决策环节，命令抵达和结果仍须另外核验。来源保留：奏折、上谕、部议、军机档或条约可证明行政动作和机构立场，不自动证明所有地区、群体或对手按其意图行动。
\eventanchor{EQ-1813-TIANLI-UPRISING-AFTERMATH}{「天理教起事者攻入紫禁城部分…}嘉庆十八年（1813年），「天理教起事者攻入紫禁城部分区域，京师发生严重安全危机」引发的善后、执行或地方回应构成可由不同材料追踪的后续节点。核心取证：《清仁宗实录》嘉庆十八年相关条；军机处档、战事方略及地方奏报；相关地区的档案、志书、契约、报刊或对方材料。该项锚定的是后续追踪问题，影响范围须按地区与群体分别判断。来源保留：战后、改革后或条约后的记忆、地方记录和官方总结常具有事后选择性；后续影响须按地点、群体和时间分层，不作整体化推断。
\eventanchor{EQ-1814-TIANLI-SUPPRESSION-PRELUDE}{「天理教起事被镇压，相关人员…}嘉庆十九年（1814年），「天理教起事被镇压，相关人员遭逮捕和处置」之前的条件、信息和争议已通过中央与地方材料显现，构成该事件可追踪的前史节点。核心取证：《清仁宗实录》嘉庆十九年相关条；宫中朱批奏折、内阁大库题本与《清史稿》相关志传；同时期地方档案、地方志、日记或对方材料应按具体地区补检。该项锚定的是前史条件与信息背景，不把条件直接写成唯一因果。来源保留：官修编年和地方材料的记录密度与立场并不相同；本行不将条件、传闻、呈报或后见解释直接等同为确定因果。
\eventanchor{EQ-1814-TIANLI-SUPPRESSION-DECISION}{围绕「天理教起事被镇压，相关…}嘉庆十九年（1814年），围绕「天理教起事被镇压，相关人员遭逮捕和处置」的命令、调度、照会或呈报形成独立的中央—地方行动文书链。核心取证：《清仁宗实录》嘉庆十九年相关条；宫中朱批奏折、内阁大库题本与《清史稿》相关志传。该项锚定的是命令或决策环节，命令抵达和结果仍须另外核验。来源保留：奏折、上谕、部议、军机档或条约可证明行政动作和机构立场，不自动证明所有地区、群体或对手按其意图行动。
\eventanchor{EQ-1814-TIANLI-SUPPRESSION-AFTERMATH}{「天理教起事被镇压，相关人员…}嘉庆十九年（1814年），「天理教起事被镇压，相关人员遭逮捕和处置」引发的善后、执行或地方回应构成可由不同材料追踪的后续节点。核心取证：《清仁宗实录》嘉庆十九年相关条；宫中朱批奏折、内阁大库题本与《清史稿》相关志传；相关地区的档案、志书、契约、报刊或对方材料。该项锚定的是后续追踪问题，影响范围须按地区与群体分别判断。来源保留：战后、改革后或条约后的记忆、地方记录和官方总结常具有事后选择性；后续影响须按地点、群体和时间分层，不作整体化推断。
\eventanchor{EQ-1815-GRAND-CANAL-REPAIRS-PRELUDE}{「漕运、河工和仓储问题在战后…}嘉庆二十年（1815年），「漕运、河工和仓储问题在战后财政压力下持续受到讨论」之前的条件、信息和争议已通过中央与地方材料显现，构成该事件可追踪的前史节点。核心取证：《清仁宗实录》嘉庆二十年相关条；户部奏销、盐政、河工或海关档案；同时期地方档案、地方志、日记或对方材料应按具体地区补检。该项锚定的是前史条件与信息背景，不把条件直接写成唯一因果。来源保留：官修编年和地方材料的记录密度与立场并不相同；本行不将条件、传闻、呈报或后见解释直接等同为确定因果。
\eventanchor{EQ-1815-GRAND-CANAL-REPAIRS-DECISION}{围绕「漕运、河工和仓储问题在…}嘉庆二十年（1815年），围绕「漕运、河工和仓储问题在战后财政压力下持续受到讨论」的命令、调度、照会或呈报形成独立的中央—地方行动文书链。核心取证：《清仁宗实录》嘉庆二十年相关条；户部奏销、盐政、河工或海关档案。该项锚定的是命令或决策环节，命令抵达和结果仍须另外核验。来源保留：奏折、上谕、部议、军机档或条约可证明行政动作和机构立场，不自动证明所有地区、群体或对手按其意图行动。
\eventanchor{EQ-1815-GRAND-CANAL-REPAIRS-AFTERMATH}{「漕运、河工和仓储问题在战后…}嘉庆二十年（1815年），「漕运、河工和仓储问题在战后财政压力下持续受到讨论」引发的善后、执行或地方回应构成可由不同材料追踪的后续节点。核心取证：《清仁宗实录》嘉庆二十年相关条；户部奏销、盐政、河工或海关档案；相关地区的档案、志书、契约、报刊或对方材料。该项锚定的是后续追踪问题，影响范围须按地区与群体分别判断。来源保留：战后、改革后或条约后的记忆、地方记录和官方总结常具有事后选择性；后续影响须按地点、群体和时间分层，不作整体化推断。
\eventanchor{EQ-1816-AMHERST-MISSION-PRELUDE}{「阿美士德使团来华，礼仪与外…}嘉庆二十一年（1816年），「阿美士德使团来华，礼仪与外交交涉未获突破」之前的条件、信息和争议已通过中央与地方材料显现，构成该事件可追踪的前史节点。核心取证：《清仁宗实录》嘉庆二十一年相关条；《筹办夷务始末》相关年卷与中外照会、条约文本；同时期地方档案、地方志、日记或对方材料应按具体地区补检。该项锚定的是前史条件与信息背景，不把条件直接写成唯一因果。来源保留：官修编年和地方材料的记录密度与立场并不相同；本行不将条件、传闻、呈报或后见解释直接等同为确定因果。
\eventanchor{EQ-1816-AMHERST-MISSION-DECISION}{围绕「阿美士德使团来华，礼仪…}嘉庆二十一年（1816年），围绕「阿美士德使团来华，礼仪与外交交涉未获突破」的命令、调度、照会或呈报形成独立的中央—地方行动文书链。核心取证：《清仁宗实录》嘉庆二十一年相关条；《筹办夷务始末》相关年卷与中外照会、条约文本。该项锚定的是命令或决策环节，命令抵达和结果仍须另外核验。来源保留：奏折、上谕、部议、军机档或条约可证明行政动作和机构立场，不自动证明所有地区、群体或对手按其意图行动。
\eventanchor{EQ-1816-AMHERST-MISSION-AFTERMATH}{「阿美士德使团来华，礼仪与外…}嘉庆二十一年（1816年），「阿美士德使团来华，礼仪与外交交涉未获突破」引发的善后、执行或地方回应构成可由不同材料追踪的后续节点。核心取证：《清仁宗实录》嘉庆二十一年相关条；《筹办夷务始末》相关年卷与中外照会、条约文本；相关地区的档案、志书、契约、报刊或对方材料。该项锚定的是后续追踪问题，影响范围须按地区与群体分别判断。来源保留：战后、改革后或条约后的记忆、地方记录和官方总结常具有事后选择性；后续影响须按地点、群体和时间分层，不作整体化推断。
\eventanchor{EQ-1818-FLOOD-RELIEF-PRELUDE}{「黄淮水患和赈务成为地方奏报…}嘉庆二十三年（1818年），「黄淮水患和赈务成为地方奏报与中央拨款的重要议题」之前的条件、信息和争议已通过中央与地方材料显现，构成该事件可追踪的前史节点。核心取证：《清仁宗实录》嘉庆二十三年相关条；地方志、督抚奏折及地方档案；同时期地方档案、地方志、日记或对方材料应按具体地区补检。该项锚定的是前史条件与信息背景，不把条件直接写成唯一因果。来源保留：官修编年和地方材料的记录密度与立场并不相同；本行不将条件、传闻、呈报或后见解释直接等同为确定因果。
\eventanchor{EQ-1818-FLOOD-RELIEF-DECISION}{围绕「黄淮水患和赈务成为地方…}嘉庆二十三年（1818年），围绕「黄淮水患和赈务成为地方奏报与中央拨款的重要议题」的命令、调度、照会或呈报形成独立的中央—地方行动文书链。核心取证：《清仁宗实录》嘉庆二十三年相关条；地方志、督抚奏折及地方档案。该项锚定的是命令或决策环节，命令抵达和结果仍须另外核验。来源保留：奏折、上谕、部议、军机档或条约可证明行政动作和机构立场，不自动证明所有地区、群体或对手按其意图行动。
\eventanchor{EQ-1818-FLOOD-RELIEF-AFTERMATH}{「黄淮水患和赈务成为地方奏报…}嘉庆二十三年（1818年），「黄淮水患和赈务成为地方奏报与中央拨款的重要议题」引发的善后、执行或地方回应构成可由不同材料追踪的后续节点。核心取证：《清仁宗实录》嘉庆二十三年相关条；地方志、督抚奏折及地方档案；相关地区的档案、志书、契约、报刊或对方材料。该项锚定的是后续追踪问题，影响范围须按地区与群体分别判断。来源保留：战后、改革后或条约后的记忆、地方记录和官方总结常具有事后选择性；后续影响须按地点、群体和时间分层，不作整体化推断。
\eventanchor{EQ-1819-OPIUM-TRADE-GROWTH-PRELUDE}{「鸦片贸易增长引起官员对白银…}嘉庆二十四年（1819年），「鸦片贸易增长引起官员对白银流动、海防和社会秩序的担忧」之前的条件、信息和争议已通过中央与地方材料显现，构成该事件可追踪的前史节点。核心取证：《清仁宗实录》嘉庆二十四年相关条；户部奏销、盐政、河工或海关档案；同时期地方档案、地方志、日记或对方材料应按具体地区补检。该项锚定的是前史条件与信息背景，不把条件直接写成唯一因果。来源保留：官修编年和地方材料的记录密度与立场并不相同；本行不将条件、传闻、呈报或后见解释直接等同为确定因果。
\eventanchor{EQ-1819-OPIUM-TRADE-GROWTH-DECISION}{围绕「鸦片贸易增长引起官员对…}嘉庆二十四年（1819年），围绕「鸦片贸易增长引起官员对白银流动、海防和社会秩序的担忧」的命令、调度、照会或呈报形成独立的中央—地方行动文书链。核心取证：《清仁宗实录》嘉庆二十四年相关条；户部奏销、盐政、河工或海关档案。该项锚定的是命令或决策环节，命令抵达和结果仍须另外核验。来源保留：奏折、上谕、部议、军机档或条约可证明行政动作和机构立场，不自动证明所有地区、群体或对手按其意图行动。
\eventanchor{EQ-1819-OPIUM-TRADE-GROWTH-AFTERMATH}{「鸦片贸易增长引起官员对白银…}嘉庆二十四年（1819年），「鸦片贸易增长引起官员对白银流动、海防和社会秩序的担忧」引发的善后、执行或地方回应构成可由不同材料追踪的后续节点。核心取证：《清仁宗实录》嘉庆二十四年相关条；户部奏销、盐政、河工或海关档案；相关地区的档案、志书、契约、报刊或对方材料。该项锚定的是后续追踪问题，影响范围须按地区与群体分别判断。来源保留：战后、改革后或条约后的记忆、地方记录和官方总结常具有事后选择性；后续影响须按地点、群体和时间分层，不作整体化推断。
\subsection{1820—1829}
\eventanchor{EQ-1820-DAO-GUANG-ACCESSION-PRELUDE}{「道光帝即位，财政、河工、旗…}嘉庆二十五年（1820年），「道光帝即位，财政、河工、旗务和沿海贸易问题进入新朝」之前的条件、信息和争议已通过中央与地方材料显现，构成该事件可追踪的前史节点。核心取证：《清仁宗实录》嘉庆二十五年相关条；宫中朱批奏折、内阁大库题本与《清史稿》相关志传；同时期地方档案、地方志、日记或对方材料应按具体地区补检。该项锚定的是前史条件与信息背景，不把条件直接写成唯一因果。来源保留：官修编年和地方材料的记录密度与立场并不相同；本行不将条件、传闻、呈报或后见解释直接等同为确定因果。
\eventanchor{EQ-1820-DAO-GUANG-ACCESSION-DECISION}{围绕「道光帝即位，财政、河工…}嘉庆二十五年（1820年），围绕「道光帝即位，财政、河工、旗务和沿海贸易问题进入新朝」的命令、调度、照会或呈报形成独立的中央—地方行动文书链。核心取证：《清仁宗实录》嘉庆二十五年相关条；宫中朱批奏折、内阁大库题本与《清史稿》相关志传。该项锚定的是命令或决策环节，命令抵达和结果仍须另外核验。来源保留：奏折、上谕、部议、军机档或条约可证明行政动作和机构立场，不自动证明所有地区、群体或对手按其意图行动。
\eventanchor{EQ-1820-DAO-GUANG-ACCESSION-AFTERMATH}{「道光帝即位，财政、河工、旗…}嘉庆二十五年（1820年），「道光帝即位，财政、河工、旗务和沿海贸易问题进入新朝」引发的善后、执行或地方回应构成可由不同材料追踪的后续节点。核心取证：《清仁宗实录》嘉庆二十五年相关条；宫中朱批奏折、内阁大库题本与《清史稿》相关志传；相关地区的档案、志书、契约、报刊或对方材料。该项锚定的是后续追踪问题，影响范围须按地区与群体分别判断。来源保留：战后、改革后或条约后的记忆、地方记录和官方总结常具有事后选择性；后续影响须按地点、群体和时间分层，不作整体化推断。
\eventanchor{EQ-1822-YELLOW-RIVER-FLOOD-PRELUDE}{「黄河水患、堤工和灾赈问题反…}道光二年（1822年），「黄河水患、堤工和灾赈问题反复进入中央议程」之前的条件、信息和争议已通过中央与地方材料显现，构成该事件可追踪的前史节点。核心取证：《清宣宗实录》道光二年相关条；地方志、督抚奏折及地方档案；同时期地方档案、地方志、日记或对方材料应按具体地区补检。该项锚定的是前史条件与信息背景，不把条件直接写成唯一因果。来源保留：官修编年和地方材料的记录密度与立场并不相同；本行不将条件、传闻、呈报或后见解释直接等同为确定因果。
\eventanchor{EQ-1822-YELLOW-RIVER-FLOOD-DECISION}{围绕「黄河水患、堤工和灾赈问…}道光二年（1822年），围绕「黄河水患、堤工和灾赈问题反复进入中央议程」的命令、调度、照会或呈报形成独立的中央—地方行动文书链。核心取证：《清宣宗实录》道光二年相关条；地方志、督抚奏折及地方档案。该项锚定的是命令或决策环节，命令抵达和结果仍须另外核验。来源保留：奏折、上谕、部议、军机档或条约可证明行政动作和机构立场，不自动证明所有地区、群体或对手按其意图行动。
\eventanchor{EQ-1822-YELLOW-RIVER-FLOOD-AFTERMATH}{「黄河水患、堤工和灾赈问题反…}道光二年（1822年），「黄河水患、堤工和灾赈问题反复进入中央议程」引发的善后、执行或地方回应构成可由不同材料追踪的后续节点。核心取证：《清宣宗实录》道光二年相关条；地方志、督抚奏折及地方档案；相关地区的档案、志书、契约、报刊或对方材料。该项锚定的是后续追踪问题，影响范围须按地区与群体分别判断。来源保留：战后、改革后或条约后的记忆、地方记录和官方总结常具有事后选择性；后续影响须按地点、群体和时间分层，不作整体化推断。
\eventanchor{EQ-1823-BANNER-DEBT-PRELUDE}{「旗人债务、典当和生计问题持…}道光三年（1823年），「旗人债务、典当和生计问题持续引起官府讨论」之前的条件、信息和争议已通过中央与地方材料显现，构成该事件可追踪的前史节点。核心取证：《清宣宗实录》道光三年相关条；地方志、督抚奏折及地方档案；同时期地方档案、地方志、日记或对方材料应按具体地区补检。该项锚定的是前史条件与信息背景，不把条件直接写成唯一因果。来源保留：官修编年和地方材料的记录密度与立场并不相同；本行不将条件、传闻、呈报或后见解释直接等同为确定因果。
\eventanchor{EQ-1823-BANNER-DEBT-DECISION}{围绕「旗人债务、典当和生计问…}道光三年（1823年），围绕「旗人债务、典当和生计问题持续引起官府讨论」的命令、调度、照会或呈报形成独立的中央—地方行动文书链。核心取证：《清宣宗实录》道光三年相关条；地方志、督抚奏折及地方档案。该项锚定的是命令或决策环节，命令抵达和结果仍须另外核验。来源保留：奏折、上谕、部议、军机档或条约可证明行政动作和机构立场，不自动证明所有地区、群体或对手按其意图行动。
\eventanchor{EQ-1823-BANNER-DEBT-AFTERMATH}{「旗人债务、典当和生计问题持…}道光三年（1823年），「旗人债务、典当和生计问题持续引起官府讨论」引发的善后、执行或地方回应构成可由不同材料追踪的后续节点。核心取证：《清宣宗实录》道光三年相关条；地方志、督抚奏折及地方档案；相关地区的档案、志书、契约、报刊或对方材料。该项锚定的是后续追踪问题，影响范围须按地区与群体分别判断。来源保留：战后、改革后或条约后的记忆、地方记录和官方总结常具有事后选择性；后续影响须按地点、群体和时间分层，不作整体化推断。
\eventanchor{EQ-1824-GRAND-CANAL-CRISIS-PRELUDE}{「运河水运和漕粮转输面临河工…}道光四年（1824年），「运河水运和漕粮转输面临河工、淤积与财政问题」之前的条件、信息和争议已通过中央与地方材料显现，构成该事件可追踪的前史节点。核心取证：《清宣宗实录》道光四年相关条；户部奏销、盐政、河工或海关档案；同时期地方档案、地方志、日记或对方材料应按具体地区补检。该项锚定的是前史条件与信息背景，不把条件直接写成唯一因果。来源保留：官修编年和地方材料的记录密度与立场并不相同；本行不将条件、传闻、呈报或后见解释直接等同为确定因果。
\eventanchor{EQ-1824-GRAND-CANAL-CRISIS-DECISION}{围绕「运河水运和漕粮转输面临…}道光四年（1824年），围绕「运河水运和漕粮转输面临河工、淤积与财政问题」的命令、调度、照会或呈报形成独立的中央—地方行动文书链。核心取证：《清宣宗实录》道光四年相关条；户部奏销、盐政、河工或海关档案。该项锚定的是命令或决策环节，命令抵达和结果仍须另外核验。来源保留：奏折、上谕、部议、军机档或条约可证明行政动作和机构立场，不自动证明所有地区、群体或对手按其意图行动。
\eventanchor{EQ-1824-GRAND-CANAL-CRISIS-AFTERMATH}{「运河水运和漕粮转输面临河工…}道光四年（1824年），「运河水运和漕粮转输面临河工、淤积与财政问题」引发的善后、执行或地方回应构成可由不同材料追踪的后续节点。核心取证：《清宣宗实录》道光四年相关条；户部奏销、盐政、河工或海关档案；相关地区的档案、志书、契约、报刊或对方材料。该项锚定的是后续追踪问题，影响范围须按地区与群体分别判断。来源保留：战后、改革后或条约后的记忆、地方记录和官方总结常具有事后选择性；后续影响须按地点、群体和时间分层，不作整体化推断。
\eventanchor{EQ-1826-JAHANGIR-KHOJA-UPRISING-PRELUDE}{「张格尔和卓进入喀什噶尔，天…}道光六年（1826年），「张格尔和卓进入喀什噶尔，天山南路爆发大规模反清动荡」之前的条件、信息和争议已通过中央与地方材料显现，构成该事件可追踪的前史节点。核心取证：《清宣宗实录》道光六年相关条；军机处档、战事方略及地方奏报；同时期地方档案、地方志、日记或对方材料应按具体地区补检。该项锚定的是前史条件与信息背景，不把条件直接写成唯一因果。来源保留：官修编年和地方材料的记录密度与立场并不相同；本行不将条件、传闻、呈报或后见解释直接等同为确定因果。
\eventanchor{EQ-1826-JAHANGIR-KHOJA-UPRISING-DECISION}{围绕「张格尔和卓进入喀什噶尔…}道光六年（1826年），围绕「张格尔和卓进入喀什噶尔，天山南路爆发大规模反清动荡」的命令、调度、照会或呈报形成独立的中央—地方行动文书链。核心取证：《清宣宗实录》道光六年相关条；军机处档、战事方略及地方奏报。该项锚定的是命令或决策环节，命令抵达和结果仍须另外核验。来源保留：奏折、上谕、部议、军机档或条约可证明行政动作和机构立场，不自动证明所有地区、群体或对手按其意图行动。
\eventanchor{EQ-1826-JAHANGIR-KHOJA-UPRISING-AFTERMATH}{「张格尔和卓进入喀什噶尔，天…}道光六年（1826年），「张格尔和卓进入喀什噶尔，天山南路爆发大规模反清动荡」引发的善后、执行或地方回应构成可由不同材料追踪的后续节点。核心取证：《清宣宗实录》道光六年相关条；军机处档、战事方略及地方奏报；相关地区的档案、志书、契约、报刊或对方材料。该项锚定的是后续追踪问题，影响范围须按地区与群体分别判断。来源保留：战后、改革后或条约后的记忆、地方记录和官方总结常具有事后选择性；后续影响须按地点、群体和时间分层，不作整体化推断。
\eventanchor{EQ-1827-KASHGAR-RECOVERY-PRELUDE}{「清军反攻并收复喀什噶尔等地…}道光七年（1827年），「清军反攻并收复喀什噶尔等地，张格尔势力受挫」之前的条件、信息和争议已通过中央与地方材料显现，构成该事件可追踪的前史节点。核心取证：《清宣宗实录》道光七年相关条；军机处档、战事方略及地方奏报；同时期地方档案、地方志、日记或对方材料应按具体地区补检。该项锚定的是前史条件与信息背景，不把条件直接写成唯一因果。来源保留：官修编年和地方材料的记录密度与立场并不相同；本行不将条件、传闻、呈报或后见解释直接等同为确定因果。
\eventanchor{EQ-1827-KASHGAR-RECOVERY-DECISION}{围绕「清军反攻并收复喀什噶尔…}道光七年（1827年），围绕「清军反攻并收复喀什噶尔等地，张格尔势力受挫」的命令、调度、照会或呈报形成独立的中央—地方行动文书链。核心取证：《清宣宗实录》道光七年相关条；军机处档、战事方略及地方奏报。该项锚定的是命令或决策环节，命令抵达和结果仍须另外核验。来源保留：奏折、上谕、部议、军机档或条约可证明行政动作和机构立场，不自动证明所有地区、群体或对手按其意图行动。
\eventanchor{EQ-1827-KASHGAR-RECOVERY-AFTERMATH}{「清军反攻并收复喀什噶尔等地…}道光七年（1827年），「清军反攻并收复喀什噶尔等地，张格尔势力受挫」引发的善后、执行或地方回应构成可由不同材料追踪的后续节点。核心取证：《清宣宗实录》道光七年相关条；军机处档、战事方略及地方奏报；相关地区的档案、志书、契约、报刊或对方材料。该项锚定的是后续追踪问题，影响范围须按地区与群体分别判断。来源保留：战后、改革后或条约后的记忆、地方记录和官方总结常具有事后选择性；后续影响须按地点、群体和时间分层，不作整体化推断。
\eventanchor{EQ-1828-JAHANGIR-CAPTURE-PRELUDE}{「张格尔被捕并押解处置，清廷…}道光八年（1828年），「张格尔被捕并押解处置，清廷加强南疆军政控制」之前的条件、信息和争议已通过中央与地方材料显现，构成该事件可追踪的前史节点。核心取证：《清宣宗实录》道光八年相关条；军机处档、理藩院档或相关方略；同时期地方档案、地方志、日记或对方材料应按具体地区补检。该项锚定的是前史条件与信息背景，不把条件直接写成唯一因果。来源保留：官修编年和地方材料的记录密度与立场并不相同；本行不将条件、传闻、呈报或后见解释直接等同为确定因果。
\eventanchor{EQ-1828-JAHANGIR-CAPTURE-DECISION}{围绕「张格尔被捕并押解处置，…}道光八年（1828年），围绕「张格尔被捕并押解处置，清廷加强南疆军政控制」的命令、调度、照会或呈报形成独立的中央—地方行动文书链。核心取证：《清宣宗实录》道光八年相关条；军机处档、理藩院档或相关方略。该项锚定的是命令或决策环节，命令抵达和结果仍须另外核验。来源保留：奏折、上谕、部议、军机档或条约可证明行政动作和机构立场，不自动证明所有地区、群体或对手按其意图行动。
\eventanchor{EQ-1828-JAHANGIR-CAPTURE-AFTERMATH}{「张格尔被捕并押解处置，清廷…}道光八年（1828年），「张格尔被捕并押解处置，清廷加强南疆军政控制」引发的善后、执行或地方回应构成可由不同材料追踪的后续节点。核心取证：《清宣宗实录》道光八年相关条；军机处档、理藩院档或相关方略；相关地区的档案、志书、契约、报刊或对方材料。该项锚定的是后续追踪问题，影响范围须按地区与群体分别判断。来源保留：战后、改革后或条约后的记忆、地方记录和官方总结常具有事后选择性；后续影响须按地点、群体和时间分层，不作整体化推断。
\subsection{1830—1839}
\eventanchor{EQ-1830-RUSSO-QING-TRADE-PRELUDE}{「恰克图等边境贸易和俄清交涉…}道光十年（1830年），「恰克图等边境贸易和俄清交涉持续影响北部边务」之前的条件、信息和争议已通过中央与地方材料显现，构成该事件可追踪的前史节点。核心取证：《清宣宗实录》道光十年相关条；《筹办夷务始末》相关年卷与中外照会、条约文本；同时期地方档案、地方志、日记或对方材料应按具体地区补检。该项锚定的是前史条件与信息背景，不把条件直接写成唯一因果。来源保留：官修编年和地方材料的记录密度与立场并不相同；本行不将条件、传闻、呈报或后见解释直接等同为确定因果。
\eventanchor{EQ-1830-RUSSO-QING-TRADE-DECISION}{围绕「恰克图等边境贸易和俄清…}道光十年（1830年），围绕「恰克图等边境贸易和俄清交涉持续影响北部边务」的命令、调度、照会或呈报形成独立的中央—地方行动文书链。核心取证：《清宣宗实录》道光十年相关条；《筹办夷务始末》相关年卷与中外照会、条约文本。该项锚定的是命令或决策环节，命令抵达和结果仍须另外核验。来源保留：奏折、上谕、部议、军机档或条约可证明行政动作和机构立场，不自动证明所有地区、群体或对手按其意图行动。
\eventanchor{EQ-1830-RUSSO-QING-TRADE-AFTERMATH}{「恰克图等边境贸易和俄清交涉…}道光十年（1830年），「恰克图等边境贸易和俄清交涉持续影响北部边务」引发的善后、执行或地方回应构成可由不同材料追踪的后续节点。核心取证：《清宣宗实录》道光十年相关条；《筹办夷务始末》相关年卷与中外照会、条约文本；相关地区的档案、志书、契约、报刊或对方材料。该项锚定的是后续追踪问题，影响范围须按地区与群体分别判断。来源保留：战后、改革后或条约后的记忆、地方记录和官方总结常具有事后选择性；后续影响须按地点、群体和时间分层，不作整体化推断。
\eventanchor{EQ-1831-CENTRAL-ASIA-CHOLERA-PRELUDE}{「疾病流行与交通网络对军队、…}道光十一年（1831年），「疾病流行与交通网络对军队、城市和商路造成压力」之前的条件、信息和争议已通过中央与地方材料显现，构成该事件可追踪的前史节点。核心取证：《清宣宗实录》道光十一年相关条；地方志、督抚奏折及地方档案；同时期地方档案、地方志、日记或对方材料应按具体地区补检。该项锚定的是前史条件与信息背景，不把条件直接写成唯一因果。来源保留：官修编年和地方材料的记录密度与立场并不相同；本行不将条件、传闻、呈报或后见解释直接等同为确定因果。
\eventanchor{EQ-1831-CENTRAL-ASIA-CHOLERA-DECISION}{围绕「疾病流行与交通网络对军…}道光十一年（1831年），围绕「疾病流行与交通网络对军队、城市和商路造成压力」的命令、调度、照会或呈报形成独立的中央—地方行动文书链。核心取证：《清宣宗实录》道光十一年相关条；地方志、督抚奏折及地方档案。该项锚定的是命令或决策环节，命令抵达和结果仍须另外核验。来源保留：奏折、上谕、部议、军机档或条约可证明行政动作和机构立场，不自动证明所有地区、群体或对手按其意图行动。
\eventanchor{EQ-1831-CENTRAL-ASIA-CHOLERA-AFTERMATH}{「疾病流行与交通网络对军队、…}道光十一年（1831年），「疾病流行与交通网络对军队、城市和商路造成压力」引发的善后、执行或地方回应构成可由不同材料追踪的后续节点。核心取证：《清宣宗实录》道光十一年相关条；地方志、督抚奏折及地方档案；相关地区的档案、志书、契约、报刊或对方材料。该项锚定的是后续追踪问题，影响范围须按地区与群体分别判断。来源保留：战后、改革后或条约后的记忆、地方记录和官方总结常具有事后选择性；后续影响须按地点、群体和时间分层，不作整体化推断。
\eventanchor{EQ-1832-COASTAL-DEFENCE-PRELUDE}{「鸦片走私增长背景下的海防、…}道光十二年（1832年），「鸦片走私增长背景下的海防、缉私和地方官责任成为争议」之前的条件、信息和争议已通过中央与地方材料显现，构成该事件可追踪的前史节点。核心取证：《清宣宗实录》道光十二年相关条；宫中朱批奏折、内阁大库题本与《清史稿》相关志传；同时期地方档案、地方志、日记或对方材料应按具体地区补检。该项锚定的是前史条件与信息背景，不把条件直接写成唯一因果。来源保留：官修编年和地方材料的记录密度与立场并不相同；本行不将条件、传闻、呈报或后见解释直接等同为确定因果。
\eventanchor{EQ-1832-COASTAL-DEFENCE-DECISION}{围绕「鸦片走私增长背景下的海…}道光十二年（1832年），围绕「鸦片走私增长背景下的海防、缉私和地方官责任成为争议」的命令、调度、照会或呈报形成独立的中央—地方行动文书链。核心取证：《清宣宗实录》道光十二年相关条；宫中朱批奏折、内阁大库题本与《清史稿》相关志传。该项锚定的是命令或决策环节，命令抵达和结果仍须另外核验。来源保留：奏折、上谕、部议、军机档或条约可证明行政动作和机构立场，不自动证明所有地区、群体或对手按其意图行动。
\eventanchor{EQ-1832-COASTAL-DEFENCE-AFTERMATH}{「鸦片走私增长背景下的海防、…}道光十二年（1832年），「鸦片走私增长背景下的海防、缉私和地方官责任成为争议」引发的善后、执行或地方回应构成可由不同材料追踪的后续节点。核心取证：《清宣宗实录》道光十二年相关条；宫中朱批奏折、内阁大库题本与《清史稿》相关志传；相关地区的档案、志书、契约、报刊或对方材料。该项锚定的是后续追踪问题，影响范围须按地区与群体分别判断。来源保留：战后、改革后或条约后的记忆、地方记录和官方总结常具有事后选择性；后续影响须按地点、群体和时间分层，不作整体化推断。
\eventanchor{EQ-1834-NAPIER-AFFAIR-PRELUDE}{「律劳卑抵广州后与两广官员发…}道光十四年（1834年），「律劳卑抵广州后与两广官员发生外交摩擦」之前的条件、信息和争议已通过中央与地方材料显现，构成该事件可追踪的前史节点。核心取证：《清宣宗实录》道光十四年相关条；《筹办夷务始末》相关年卷与中外照会、条约文本；同时期地方档案、地方志、日记或对方材料应按具体地区补检。该项锚定的是前史条件与信息背景，不把条件直接写成唯一因果。来源保留：官修编年和地方材料的记录密度与立场并不相同；本行不将条件、传闻、呈报或后见解释直接等同为确定因果。
\eventanchor{EQ-1834-NAPIER-AFFAIR-DECISION}{围绕「律劳卑抵广州后与两广官…}道光十四年（1834年），围绕「律劳卑抵广州后与两广官员发生外交摩擦」的命令、调度、照会或呈报形成独立的中央—地方行动文书链。核心取证：《清宣宗实录》道光十四年相关条；《筹办夷务始末》相关年卷与中外照会、条约文本。该项锚定的是命令或决策环节，命令抵达和结果仍须另外核验。来源保留：奏折、上谕、部议、军机档或条约可证明行政动作和机构立场，不自动证明所有地区、群体或对手按其意图行动。
\eventanchor{EQ-1834-NAPIER-AFFAIR-AFTERMATH}{「律劳卑抵广州后与两广官员发…}道光十四年（1834年），「律劳卑抵广州后与两广官员发生外交摩擦」引发的善后、执行或地方回应构成可由不同材料追踪的后续节点。核心取证：《清宣宗实录》道光十四年相关条；《筹办夷务始末》相关年卷与中外照会、条约文本；相关地区的档案、志书、契约、报刊或对方材料。该项锚定的是后续追踪问题，影响范围须按地区与群体分别判断。来源保留：战后、改革后或条约后的记忆、地方记录和官方总结常具有事后选择性；后续影响须按地点、群体和时间分层，不作整体化推断。
\eventanchor{EQ-1835-GUANGZHOU-TRADE-REGULATION-PRELUDE}{「广州贸易、行商债务和外商管…}道光十五年（1835年），「广州贸易、行商债务和外商管理继续在地方与中央之间协商」之前的条件、信息和争议已通过中央与地方材料显现，构成该事件可追踪的前史节点。核心取证：《清宣宗实录》道光十五年相关条；户部奏销、盐政、河工或海关档案；同时期地方档案、地方志、日记或对方材料应按具体地区补检。该项锚定的是前史条件与信息背景，不把条件直接写成唯一因果。来源保留：官修编年和地方材料的记录密度与立场并不相同；本行不将条件、传闻、呈报或后见解释直接等同为确定因果。
\eventanchor{EQ-1835-GUANGZHOU-TRADE-REGULATION-DECISION}{围绕「广州贸易、行商债务和外…}道光十五年（1835年），围绕「广州贸易、行商债务和外商管理继续在地方与中央之间协商」的命令、调度、照会或呈报形成独立的中央—地方行动文书链。核心取证：《清宣宗实录》道光十五年相关条；户部奏销、盐政、河工或海关档案。该项锚定的是命令或决策环节，命令抵达和结果仍须另外核验。来源保留：奏折、上谕、部议、军机档或条约可证明行政动作和机构立场，不自动证明所有地区、群体或对手按其意图行动。
\eventanchor{EQ-1835-GUANGZHOU-TRADE-REGULATION-AFTERMATH}{「广州贸易、行商债务和外商管…}道光十五年（1835年），「广州贸易、行商债务和外商管理继续在地方与中央之间协商」引发的善后、执行或地方回应构成可由不同材料追踪的后续节点。核心取证：《清宣宗实录》道光十五年相关条；户部奏销、盐政、河工或海关档案；相关地区的档案、志书、契约、报刊或对方材料。该项锚定的是后续追踪问题，影响范围须按地区与群体分别判断。来源保留：战后、改革后或条约后的记忆、地方记录和官方总结常具有事后选择性；后续影响须按地点、群体和时间分层，不作整体化推断。
\eventanchor{EQ-1836-OPIUM-LEGALIZATION-DEBATE-PRELUDE}{「许乃济等围绕鸦片弛禁、征税…}道光十六年（1836年），「许乃济等围绕鸦片弛禁、征税与禁绝提出政策论争」之前的条件、信息和争议已通过中央与地方材料显现，构成该事件可追踪的前史节点。核心取证：《清宣宗实录》道光十六年相关条；宫中朱批奏折、内阁大库题本与《清史稿》相关志传；同时期地方档案、地方志、日记或对方材料应按具体地区补检。该项锚定的是前史条件与信息背景，不把条件直接写成唯一因果。来源保留：官修编年和地方材料的记录密度与立场并不相同；本行不将条件、传闻、呈报或后见解释直接等同为确定因果。
\eventanchor{EQ-1836-OPIUM-LEGALIZATION-DEBATE-DECISION}{围绕「许乃济等围绕鸦片弛禁、…}道光十六年（1836年），围绕「许乃济等围绕鸦片弛禁、征税与禁绝提出政策论争」的命令、调度、照会或呈报形成独立的中央—地方行动文书链。核心取证：《清宣宗实录》道光十六年相关条；宫中朱批奏折、内阁大库题本与《清史稿》相关志传。该项锚定的是命令或决策环节，命令抵达和结果仍须另外核验。来源保留：奏折、上谕、部议、军机档或条约可证明行政动作和机构立场，不自动证明所有地区、群体或对手按其意图行动。
\eventanchor{EQ-1836-OPIUM-LEGALIZATION-DEBATE-AFTERMATH}{「许乃济等围绕鸦片弛禁、征税…}道光十六年（1836年），「许乃济等围绕鸦片弛禁、征税与禁绝提出政策论争」引发的善后、执行或地方回应构成可由不同材料追踪的后续节点。核心取证：《清宣宗实录》道光十六年相关条；宫中朱批奏折、内阁大库题本与《清史稿》相关志传；相关地区的档案、志书、契约、报刊或对方材料。该项锚定的是后续追踪问题，影响范围须按地区与群体分别判断。来源保留：战后、改革后或条约后的记忆、地方记录和官方总结常具有事后选择性；后续影响须按地点、群体和时间分层，不作整体化推断。
\eventanchor{EQ-1837-CANTON-SMUGGLING-PRELUDE}{「广东缉私与鸦片走私之间的冲…}道光十七年（1837年），「广东缉私与鸦片走私之间的冲突持续升级」之前的条件、信息和争议已通过中央与地方材料显现，构成该事件可追踪的前史节点。核心取证：《清宣宗实录》道光十七年相关条；地方志、督抚奏折及地方档案；同时期地方档案、地方志、日记或对方材料应按具体地区补检。该项锚定的是前史条件与信息背景，不把条件直接写成唯一因果。来源保留：官修编年和地方材料的记录密度与立场并不相同；本行不将条件、传闻、呈报或后见解释直接等同为确定因果。
\eventanchor{EQ-1837-CANTON-SMUGGLING-DECISION}{围绕「广东缉私与鸦片走私之间…}道光十七年（1837年），围绕「广东缉私与鸦片走私之间的冲突持续升级」的命令、调度、照会或呈报形成独立的中央—地方行动文书链。核心取证：《清宣宗实录》道光十七年相关条；地方志、督抚奏折及地方档案。该项锚定的是命令或决策环节，命令抵达和结果仍须另外核验。来源保留：奏折、上谕、部议、军机档或条约可证明行政动作和机构立场，不自动证明所有地区、群体或对手按其意图行动。
\eventanchor{EQ-1837-CANTON-SMUGGLING-AFTERMATH}{「广东缉私与鸦片走私之间的冲…}道光十七年（1837年），「广东缉私与鸦片走私之间的冲突持续升级」引发的善后、执行或地方回应构成可由不同材料追踪的后续节点。核心取证：《清宣宗实录》道光十七年相关条；地方志、督抚奏折及地方档案；相关地区的档案、志书、契约、报刊或对方材料。该项锚定的是后续追踪问题，影响范围须按地区与群体分别判断。来源保留：战后、改革后或条约后的记忆、地方记录和官方总结常具有事后选择性；后续影响须按地点、群体和时间分层，不作整体化推断。
\eventanchor{EQ-1839-HUMEN-OPIUM-DESTRUCTION-PRELUDE}{「林则徐在虎门销毁收缴鸦片，…}道光十九年（1839年），「林则徐在虎门销毁收缴鸦片，禁烟冲突公开化」之前的条件、信息和争议已通过中央与地方材料显现，构成该事件可追踪的前史节点。核心取证：《清宣宗实录》道光十九年相关条；《筹办夷务始末》相关年卷与中外照会、条约文本；同时期地方档案、地方志、日记或对方材料应按具体地区补检。该项锚定的是前史条件与信息背景，不把条件直接写成唯一因果。来源保留：官修编年和地方材料的记录密度与立场并不相同；本行不将条件、传闻、呈报或后见解释直接等同为确定因果。
\eventanchor{EQ-1839-HUMEN-OPIUM-DESTRUCTION-DECISION}{围绕「林则徐在虎门销毁收缴鸦…}道光十九年（1839年），围绕「林则徐在虎门销毁收缴鸦片，禁烟冲突公开化」的命令、调度、照会或呈报形成独立的中央—地方行动文书链。核心取证：《清宣宗实录》道光十九年相关条；《筹办夷务始末》相关年卷与中外照会、条约文本。该项锚定的是命令或决策环节，命令抵达和结果仍须另外核验。来源保留：奏折、上谕、部议、军机档或条约可证明行政动作和机构立场，不自动证明所有地区、群体或对手按其意图行动。
\eventanchor{EQ-1839-HUMEN-OPIUM-DESTRUCTION-AFTERMATH}{「林则徐在虎门销毁收缴鸦片，…}道光十九年（1839年），「林则徐在虎门销毁收缴鸦片，禁烟冲突公开化」引发的善后、执行或地方回应构成可由不同材料追踪的后续节点。核心取证：《清宣宗实录》道光十九年相关条；《筹办夷务始末》相关年卷与中外照会、条约文本；相关地区的档案、志书、契约、报刊或对方材料。该项锚定的是后续追踪问题，影响范围须按地区与群体分别判断。来源保留：战后、改革后或条约后的记忆、地方记录和官方总结常具有事后选择性；后续影响须按地点、群体和时间分层，不作整体化推断。
\subsection{1840—1849}
\eventanchor{EQ-1840-FIRST-OPIUM-WAR-PRELUDE}{「英国舰队来华，第一次鸦片战…}道光二十年（1840年），「英国舰队来华，第一次鸦片战争爆发」之前的条件、信息和争议已通过中央与地方材料显现，构成该事件可追踪的前史节点。核心取证：《清宣宗实录》道光二十年相关条；军机处档、战事方略及地方奏报；同时期地方档案、地方志、日记或对方材料应按具体地区补检。该项锚定的是前史条件与信息背景，不把条件直接写成唯一因果。来源保留：官修编年和地方材料的记录密度与立场并不相同；本行不将条件、传闻、呈报或后见解释直接等同为确定因果。
\eventanchor{EQ-1840-FIRST-OPIUM-WAR-DECISION}{围绕「英国舰队来华，第一次鸦…}道光二十年（1840年），围绕「英国舰队来华，第一次鸦片战争爆发」的命令、调度、照会或呈报形成独立的中央—地方行动文书链。核心取证：《清宣宗实录》道光二十年相关条；军机处档、战事方略及地方奏报。该项锚定的是命令或决策环节，命令抵达和结果仍须另外核验。来源保留：奏折、上谕、部议、军机档或条约可证明行政动作和机构立场，不自动证明所有地区、群体或对手按其意图行动。
\eventanchor{EQ-1840-FIRST-OPIUM-WAR-AFTERMATH}{「英国舰队来华，第一次鸦片战…}道光二十年（1840年），「英国舰队来华，第一次鸦片战争爆发」引发的善后、执行或地方回应构成可由不同材料追踪的后续节点。核心取证：《清宣宗实录》道光二十年相关条；军机处档、战事方略及地方奏报；相关地区的档案、志书、契约、报刊或对方材料。该项锚定的是后续追踪问题，影响范围须按地区与群体分别判断。来源保留：战后、改革后或条约后的记忆、地方记录和官方总结常具有事后选择性；后续影响须按地点、群体和时间分层，不作整体化推断。
\eventanchor{EQ-1841-CHUANBI-NEGOTIATIONS-PRELUDE}{「穿鼻草约等谈判和沿海战事并…}道光二十一年（1841年），「穿鼻草约等谈判和沿海战事并行，清廷内部对和战意见分歧」之前的条件、信息和争议已通过中央与地方材料显现，构成该事件可追踪的前史节点。核心取证：《清宣宗实录》道光二十一年相关条；《筹办夷务始末》相关年卷与中外照会、条约文本；同时期地方档案、地方志、日记或对方材料应按具体地区补检。该项锚定的是前史条件与信息背景，不把条件直接写成唯一因果。来源保留：官修编年和地方材料的记录密度与立场并不相同；本行不将条件、传闻、呈报或后见解释直接等同为确定因果。
\eventanchor{EQ-1841-CHUANBI-NEGOTIATIONS-DECISION}{围绕「穿鼻草约等谈判和沿海战…}道光二十一年（1841年），围绕「穿鼻草约等谈判和沿海战事并行，清廷内部对和战意见分歧」的命令、调度、照会或呈报形成独立的中央—地方行动文书链。核心取证：《清宣宗实录》道光二十一年相关条；《筹办夷务始末》相关年卷与中外照会、条约文本。该项锚定的是命令或决策环节，命令抵达和结果仍须另外核验。来源保留：奏折、上谕、部议、军机档或条约可证明行政动作和机构立场，不自动证明所有地区、群体或对手按其意图行动。
\eventanchor{EQ-1841-CHUANBI-NEGOTIATIONS-AFTERMATH}{「穿鼻草约等谈判和沿海战事并…}道光二十一年（1841年），「穿鼻草约等谈判和沿海战事并行，清廷内部对和战意见分歧」引发的善后、执行或地方回应构成可由不同材料追踪的后续节点。核心取证：《清宣宗实录》道光二十一年相关条；《筹办夷务始末》相关年卷与中外照会、条约文本；相关地区的档案、志书、契约、报刊或对方材料。该项锚定的是后续追踪问题，影响范围须按地区与群体分别判断。来源保留：战后、改革后或条约后的记忆、地方记录和官方总结常具有事后选择性；后续影响须按地点、群体和时间分层，不作整体化推断。
\eventanchor{EQ-1842-NANJING-TREATY-PRELUDE}{「《南京条约》签订，清廷割让…}道光二十二年（1842年），「《南京条约》签订，清廷割让香港并开放五口通商、支付赔款」之前的条件、信息和争议已通过中央与地方材料显现，构成该事件可追踪的前史节点。核心取证：《清宣宗实录》道光二十二年相关条；《筹办夷务始末》相关年卷与中外照会、条约文本；同时期地方档案、地方志、日记或对方材料应按具体地区补检。该项锚定的是前史条件与信息背景，不把条件直接写成唯一因果。来源保留：官修编年和地方材料的记录密度与立场并不相同；本行不将条件、传闻、呈报或后见解释直接等同为确定因果。
\eventanchor{EQ-1842-NANJING-TREATY-DECISION}{围绕「《南京条约》签订，清廷…}道光二十二年（1842年），围绕「《南京条约》签订，清廷割让香港并开放五口通商、支付赔款」的命令、调度、照会或呈报形成独立的中央—地方行动文书链。核心取证：《清宣宗实录》道光二十二年相关条；《筹办夷务始末》相关年卷与中外照会、条约文本。该项锚定的是命令或决策环节，命令抵达和结果仍须另外核验。来源保留：奏折、上谕、部议、军机档或条约可证明行政动作和机构立场，不自动证明所有地区、群体或对手按其意图行动。
\eventanchor{EQ-1842-NANJING-TREATY-AFTERMATH}{「《南京条约》签订，清廷割让…}道光二十二年（1842年），「《南京条约》签订，清廷割让香港并开放五口通商、支付赔款」引发的善后、执行或地方回应构成可由不同材料追踪的后续节点。核心取证：《清宣宗实录》道光二十二年相关条；《筹办夷务始末》相关年卷与中外照会、条约文本；相关地区的档案、志书、契约、报刊或对方材料。该项锚定的是后续追踪问题，影响范围须按地区与群体分别判断。来源保留：战后、改革后或条约后的记忆、地方记录和官方总结常具有事后选择性；后续影响须按地点、群体和时间分层，不作整体化推断。
\eventanchor{EQ-1843-BOGUE-TREATY-PRELUDE}{「《虎门条约》补充领事裁判、…}道光二十三年（1843年），「《虎门条约》补充领事裁判、协定关税等安排」之前的条件、信息和争议已通过中央与地方材料显现，构成该事件可追踪的前史节点。核心取证：《清宣宗实录》道光二十三年相关条；《筹办夷务始末》相关年卷与中外照会、条约文本；同时期地方档案、地方志、日记或对方材料应按具体地区补检。该项锚定的是前史条件与信息背景，不把条件直接写成唯一因果。来源保留：官修编年和地方材料的记录密度与立场并不相同；本行不将条件、传闻、呈报或后见解释直接等同为确定因果。
\eventanchor{EQ-1843-BOGUE-TREATY-DECISION}{围绕「《虎门条约》补充领事裁…}道光二十三年（1843年），围绕「《虎门条约》补充领事裁判、协定关税等安排」的命令、调度、照会或呈报形成独立的中央—地方行动文书链。核心取证：《清宣宗实录》道光二十三年相关条；《筹办夷务始末》相关年卷与中外照会、条约文本。该项锚定的是命令或决策环节，命令抵达和结果仍须另外核验。来源保留：奏折、上谕、部议、军机档或条约可证明行政动作和机构立场，不自动证明所有地区、群体或对手按其意图行动。
\eventanchor{EQ-1843-BOGUE-TREATY-AFTERMATH}{「《虎门条约》补充领事裁判、…}道光二十三年（1843年），「《虎门条约》补充领事裁判、协定关税等安排」引发的善后、执行或地方回应构成可由不同材料追踪的后续节点。核心取证：《清宣宗实录》道光二十三年相关条；《筹办夷务始末》相关年卷与中外照会、条约文本；相关地区的档案、志书、契约、报刊或对方材料。该项锚定的是后续追踪问题，影响范围须按地区与群体分别判断。来源保留：战后、改革后或条约后的记忆、地方记录和官方总结常具有事后选择性；后续影响须按地点、群体和时间分层，不作整体化推断。
\eventanchor{EQ-1844-WANGXIA-WHAMPOA-TREATIES-PRELUDE}{「《望厦条约》《黄埔条约》分…}道光二十四年（1844年），「《望厦条约》《黄埔条约》分别确立美法在华通商与相关权利」之前的条件、信息和争议已通过中央与地方材料显现，构成该事件可追踪的前史节点。核心取证：《清宣宗实录》道光二十四年相关条；《筹办夷务始末》相关年卷与中外照会、条约文本；同时期地方档案、地方志、日记或对方材料应按具体地区补检。该项锚定的是前史条件与信息背景，不把条件直接写成唯一因果。来源保留：官修编年和地方材料的记录密度与立场并不相同；本行不将条件、传闻、呈报或后见解释直接等同为确定因果。
\eventanchor{EQ-1844-WANGXIA-WHAMPOA-TREATIES-DECISION}{围绕「《望厦条约》《黄埔条约…}道光二十四年（1844年），围绕「《望厦条约》《黄埔条约》分别确立美法在华通商与相关权利」的命令、调度、照会或呈报形成独立的中央—地方行动文书链。核心取证：《清宣宗实录》道光二十四年相关条；《筹办夷务始末》相关年卷与中外照会、条约文本。该项锚定的是命令或决策环节，命令抵达和结果仍须另外核验。来源保留：奏折、上谕、部议、军机档或条约可证明行政动作和机构立场，不自动证明所有地区、群体或对手按其意图行动。
\eventanchor{EQ-1844-WANGXIA-WHAMPOA-TREATIES-AFTERMATH}{「《望厦条约》《黄埔条约》分…}道光二十四年（1844年），「《望厦条约》《黄埔条约》分别确立美法在华通商与相关权利」引发的善后、执行或地方回应构成可由不同材料追踪的后续节点。核心取证：《清宣宗实录》道光二十四年相关条；《筹办夷务始末》相关年卷与中外照会、条约文本；相关地区的档案、志书、契约、报刊或对方材料。该项锚定的是后续追踪问题，影响范围须按地区与群体分别判断。来源保留：战后、改革后或条约后的记忆、地方记录和官方总结常具有事后选择性；后续影响须按地点、群体和时间分层，不作整体化推断。
\eventanchor{EQ-1846-MISSIONARY-RETURN-PRELUDE}{「条约环境下传教活动与地方教…}道光二十六年（1846年），「条约环境下传教活动与地方教案问题逐渐增多」之前的条件、信息和争议已通过中央与地方材料显现，构成该事件可追踪的前史节点。核心取证：《清宣宗实录》道光二十六年相关条；地方志、督抚奏折及地方档案；同时期地方档案、地方志、日记或对方材料应按具体地区补检。该项锚定的是前史条件与信息背景，不把条件直接写成唯一因果。来源保留：官修编年和地方材料的记录密度与立场并不相同；本行不将条件、传闻、呈报或后见解释直接等同为确定因果。
\eventanchor{EQ-1846-MISSIONARY-RETURN-DECISION}{围绕「条约环境下传教活动与地…}道光二十六年（1846年），围绕「条约环境下传教活动与地方教案问题逐渐增多」的命令、调度、照会或呈报形成独立的中央—地方行动文书链。核心取证：《清宣宗实录》道光二十六年相关条；地方志、督抚奏折及地方档案。该项锚定的是命令或决策环节，命令抵达和结果仍须另外核验。来源保留：奏折、上谕、部议、军机档或条约可证明行政动作和机构立场，不自动证明所有地区、群体或对手按其意图行动。
\eventanchor{EQ-1846-MISSIONARY-RETURN-AFTERMATH}{「条约环境下传教活动与地方教…}道光二十六年（1846年），「条约环境下传教活动与地方教案问题逐渐增多」引发的善后、执行或地方回应构成可由不同材料追踪的后续节点。核心取证：《清宣宗实录》道光二十六年相关条；地方志、督抚奏折及地方档案；相关地区的档案、志书、契约、报刊或对方材料。该项锚定的是后续追踪问题，影响范围须按地区与群体分别判断。来源保留：战后、改革后或条约后的记忆、地方记录和官方总结常具有事后选择性；后续影响须按地点、群体和时间分层，不作整体化推断。
\eventanchor{EQ-1847-GUANGZHOU-CONFLICT-PRELUDE}{「广州周边的中英冲突显示条约…}道光二十七年（1847年），「广州周边的中英冲突显示条约执行和城市主权问题尚未解决」之前的条件、信息和争议已通过中央与地方材料显现，构成该事件可追踪的前史节点。核心取证：《清宣宗实录》道光二十七年相关条；《筹办夷务始末》相关年卷与中外照会、条约文本；同时期地方档案、地方志、日记或对方材料应按具体地区补检。该项锚定的是前史条件与信息背景，不把条件直接写成唯一因果。来源保留：官修编年和地方材料的记录密度与立场并不相同；本行不将条件、传闻、呈报或后见解释直接等同为确定因果。
\eventanchor{EQ-1847-GUANGZHOU-CONFLICT-DECISION}{围绕「广州周边的中英冲突显示…}道光二十七年（1847年），围绕「广州周边的中英冲突显示条约执行和城市主权问题尚未解决」的命令、调度、照会或呈报形成独立的中央—地方行动文书链。核心取证：《清宣宗实录》道光二十七年相关条；《筹办夷务始末》相关年卷与中外照会、条约文本。该项锚定的是命令或决策环节，命令抵达和结果仍须另外核验。来源保留：奏折、上谕、部议、军机档或条约可证明行政动作和机构立场，不自动证明所有地区、群体或对手按其意图行动。
\eventanchor{EQ-1847-GUANGZHOU-CONFLICT-AFTERMATH}{「广州周边的中英冲突显示条约…}道光二十七年（1847年），「广州周边的中英冲突显示条约执行和城市主权问题尚未解决」引发的善后、执行或地方回应构成可由不同材料追踪的后续节点。核心取证：《清宣宗实录》道光二十七年相关条；《筹办夷务始末》相关年卷与中外照会、条约文本；相关地区的档案、志书、契约、报刊或对方材料。该项锚定的是后续追踪问题，影响范围须按地区与群体分别判断。来源保留：战后、改革后或条约后的记忆、地方记录和官方总结常具有事后选择性；后续影响须按地点、群体和时间分层，不作整体化推断。
\eventanchor{EQ-1848-YANGTZE-TRADE-PRELUDE}{「通商口岸贸易网络向长江和沿…}道光二十八年（1848年），「通商口岸贸易网络向长江和沿海延伸，地方商人开始适应新制度」之前的条件、信息和争议已通过中央与地方材料显现，构成该事件可追踪的前史节点。核心取证：《清宣宗实录》道光二十八年相关条；户部奏销、盐政、河工或海关档案；同时期地方档案、地方志、日记或对方材料应按具体地区补检。该项锚定的是前史条件与信息背景，不把条件直接写成唯一因果。来源保留：官修编年和地方材料的记录密度与立场并不相同；本行不将条件、传闻、呈报或后见解释直接等同为确定因果。
\eventanchor{EQ-1848-YANGTZE-TRADE-DECISION}{围绕「通商口岸贸易网络向长江…}道光二十八年（1848年），围绕「通商口岸贸易网络向长江和沿海延伸，地方商人开始适应新制度」的命令、调度、照会或呈报形成独立的中央—地方行动文书链。核心取证：《清宣宗实录》道光二十八年相关条；户部奏销、盐政、河工或海关档案。该项锚定的是命令或决策环节，命令抵达和结果仍须另外核验。来源保留：奏折、上谕、部议、军机档或条约可证明行政动作和机构立场，不自动证明所有地区、群体或对手按其意图行动。
\eventanchor{EQ-1848-YANGTZE-TRADE-AFTERMATH}{「通商口岸贸易网络向长江和沿…}道光二十八年（1848年），「通商口岸贸易网络向长江和沿海延伸，地方商人开始适应新制度」引发的善后、执行或地方回应构成可由不同材料追踪的后续节点。核心取证：《清宣宗实录》道光二十八年相关条；户部奏销、盐政、河工或海关档案；相关地区的档案、志书、契约、报刊或对方材料。该项锚定的是后续追踪问题，影响范围须按地区与群体分别判断。来源保留：战后、改革后或条约后的记忆、地方记录和官方总结常具有事后选择性；后续影响须按地点、群体和时间分层，不作整体化推断。
\subsection{1850—1859}
\eventanchor{EQ-1850-XIANFENG-ACCESSION-PRELUDE}{「咸丰帝即位，财政紧张、内乱…}道光三十年（1850年），「咸丰帝即位，财政紧张、内乱和对外压力同时加重」之前的条件、信息和争议已通过中央与地方材料显现，构成该事件可追踪的前史节点。核心取证：《清宣宗实录》道光三十年相关条；宫中朱批奏折、内阁大库题本与《清史稿》相关志传；同时期地方档案、地方志、日记或对方材料应按具体地区补检。该项锚定的是前史条件与信息背景，不把条件直接写成唯一因果。来源保留：官修编年和地方材料的记录密度与立场并不相同；本行不将条件、传闻、呈报或后见解释直接等同为确定因果。
\eventanchor{EQ-1850-XIANFENG-ACCESSION-DECISION}{围绕「咸丰帝即位，财政紧张、…}道光三十年（1850年），围绕「咸丰帝即位，财政紧张、内乱和对外压力同时加重」的命令、调度、照会或呈报形成独立的中央—地方行动文书链。核心取证：《清宣宗实录》道光三十年相关条；宫中朱批奏折、内阁大库题本与《清史稿》相关志传。该项锚定的是命令或决策环节，命令抵达和结果仍须另外核验。来源保留：奏折、上谕、部议、军机档或条约可证明行政动作和机构立场，不自动证明所有地区、群体或对手按其意图行动。
\eventanchor{EQ-1850-XIANFENG-ACCESSION-AFTERMATH}{「咸丰帝即位，财政紧张、内乱…}道光三十年（1850年），「咸丰帝即位，财政紧张、内乱和对外压力同时加重」引发的善后、执行或地方回应构成可由不同材料追踪的后续节点。核心取证：《清宣宗实录》道光三十年相关条；宫中朱批奏折、内阁大库题本与《清史稿》相关志传；相关地区的档案、志书、契约、报刊或对方材料。该项锚定的是后续追踪问题，影响范围须按地区与群体分别判断。来源保留：战后、改革后或条约后的记忆、地方记录和官方总结常具有事后选择性；后续影响须按地点、群体和时间分层，不作整体化推断。
